\tituloI{METODOLOGÍA}

\tituloII{Método, tipo y alcance de la investigación}

\tituloIII{Método de investigación}
\parrafoIII{
    Se empleará el \textbf{método analítico-sintético}. Este proceso permitirá descomponer el fenómeno del uso de la IA en sus indicadores básicos (análisis) para luego integrar los hallazgos en una interpretación global sobre la formación académica en la sede Cusco (síntesis) \cite{arias2012}. Asimismo, se aplicará el método deductivo para transitar de las teorías generales de la adopción tecnológica hacia la realidad particular del estudiante de ingeniería.
}

\tituloIII{Tipo de investigación}
\parrafoIII{
    La investigación es de \textbf{tipo aplicada}. Según la finalidad, este estudio busca utilizar los conocimientos teóricos sobre inteligencia artificial para diagnosticar y proponer mejoras en las estrategias de aprendizaje de la Escuela Profesional de Ingeniería de Sistemas e Informática \cite{arias2012}.
}

\tituloIII{Alcance de la investigación}
\parrafoIII{
    El estudio posee un \textbf{alcance descriptivo}. Su propósito es especificar las propiedades y características importantes de la población estudiantil respecto a la variable en estudio, midiendo cada dimensión de manera independiente para describir el fenómeno tal como se manifiesta en su contexto natural \cite{hernandez2014}.
}

\tituloII{Materiales y métodos}

\tituloIII{Materiales}
\parrafoIII{
    Para el desarrollo de la investigación se utilizarán los siguientes recursos técnicos y lógicos:
}
\begin{listaIII}
    \item \textbf{Hardware:} Computadora personal con procesador de arquitectura x86\_64 y periféricos de entrada para el procesamiento de datos.
    \item \textbf{Software de Procesamiento:} Sistema operativo basado en el kernel Linux (distribución Arch Linux) para garantizar la estabilidad del entorno de trabajo.
    \item \textbf{Herramientas de Análisis:} Suite estadística GNU R o Microsoft Excel para la tabulación de datos y el motor de composición de textos LaTeX (TeX Live) para la redacción del informe final.
    \item \textbf{Instrumento:} Cuestionario digital distribuido mediante formularios web.
\end{listaIII}

\tituloIII{Métodos}
\parrafoIII{
    El procedimiento se dividirá en tres fases: 1) Fase de gabinete, donde se diseñará el cuestionario basado en la operacionalización de variables; 2) Fase de campo, donde se aplicará el instrumento a la muestra de estudiantes en la sede San Jerónimo; y 3) Fase de procesamiento, donde se realizará el análisis descriptivo mediante estadística de frecuencias y porcentajes.
}